% ===== §3 The Generation of Perception: Several Schools =====
\section{The Generation of Perception: Several Schools}
\label{sec:generation}

\noindent
The first spine of this paper holds that aesthetic perception generates a world rather than registering one. The claim is easy to assert and easy to reduce to a slogan, and it earns its place only if it can be shown to be the convergent finding of several traditions that arrived at it by different routes and in different vocabularies. This section assembles those traditions. It takes up in turn the phenomenology of perception, the predictive account of the perceiving brain, and the psychoanalytic account of how desire organises the visual field, and it shows that each, pressed on the question of what perceiving is, returns a version of the same answer: that perception is a construction of the world under a prior frame, and not the passive receipt of a world already made. It then meets the objection that would dissolve the whole enterprise, the denial that there is any distinctively aesthetic perception at all, and closes with the contemporary aesthetics of atmosphere, which locates the beautiful in exactly the relational between that this paper's ontology requires. The traditions are not merged here. Their convergence is the business of a later chapter (\xref{sec:convergence}); the business of this one is to let each speak.

\subsection{The phenomenological genealogy}
\label{sub:gen-phenomenology}

Phenomenology reached the constructive character of perception first and states it most directly. For Merleau-Ponty, perception is not the reception of sense-data by a mind that then interprets them; it is the activity of a body already geared into a world, a body whose \emph{schema} is a prior organisation of the perceptual field in advance of any particular thing perceived within it \citep{merleau-ponty1962}. To perceive is to bring a bodily hold on the world to bear on what appears, so that what appears is always already articulated by the hold. There is no moment of pure reception prior to this articulation; the articulation is perception. Husserl, working the same ground from another side, gave the point its most exact name in the analyses of passive synthesis: before any active judgement, the perceptual field is already synthesised, already organised into unities and saliences and continuities by a passivity that has done its constructive work beneath and prior to the perceiver's notice \citep{husserl2001}. What arrives at active attention has already been shaped. The everyday is the natural home of this passivity, since it is precisely in the unattended flow of ordinary perception that passive synthesis does the greater part of its constructive labour, delivering a world already organised to a perceiver who takes that organisation for the world itself. Phenomenology's contribution to the first spine is therefore this: that the givenness of the perceived world is the achievement of a constructive activity, bodily and passive, which perception mistakes for a mere finding because the construction is not itself perceived.

\subsection{The cognitive-scientific genealogy}
\label{sub:gen-cognitive}

The predictive account of the brain reaches the same conclusion through the study of the perceiving nervous system, and it is worth marking that it was not built to vindicate phenomenology and vindicates it anyway. On this account, perception is not the upward flow of sensory information into a receiving mind but a process in which the perceptual system continually generates predictions of its sensory input from a prior model of the world and corrects that model only by the residual error its predictions fail to cancel \citep{friston2010}. What is perceived is, in the first instance, what the system predicted; sensation enters as the correction of prediction rather than as the origin of perception. The frame under which the phenomenologists said perception constructs its object is, in this idiom, the generative model the brain brings to its input, and the constructive priority they described is here the literal direction of the dominant signal, which runs from model to world before it runs from world to model. The paper treats this correspondence as a substantive convergence and not a loose analogy, and it defers to the chapter on convergence the argument that the two are versions of one finding, and to the chapter on plasticity (\xref{sec:plasticity}) the neural body of the claim. What the first spine draws from cognitive science here is only the philosophical point: that the best current theory of the perceiving brain describes perception as the generation of a predicted world, corrected by sensation, rather than the assembly of a world out of sensation.

\subsection{The psychoanalytic genealogy}
\label{sub:gen-psychoanalysis}

Phenomenology and cognitive science establish that perception constructs under a prior frame; psychoanalysis asks what organises the frame, and answers with desire. Its contribution to the first spine is the one the other two cannot supply, an account of why the constructed world is constructed \emph{thus}, salient here and blind there, and its answer is that the perceptual field is organised around desire and its object. Lacan's analysis of the visual field turns on the claim that vision is structured by a gaze that does not belong to the seeing subject, a gaze that functions as \emph{objet a}, the object-cause of desire, around which the scopic field is arranged \citep{lacan1978}. What a perceiver sees, and the point from which the seen is organised into a scene, is set by a desire that precedes and exceeds the perceiver's intention; the fantasy frame decides, beneath consciousness, what will show up as visible and what will remain unregistered. This is not a mysticism. It is the specification, at the level of desire, of the prior frame the other two traditions described formally: the generative model is not neutral, the bodily hold is not disinterested, and what organises their selectivity is, in the cases this paper cares about, the structure of a desire. The consequence for an aesthetics of the everyday is severe and clarifying. The beauty a perceiver finds in an ordinary scene is, in part, the beauty their desire was structured to find; the light at the window is seen as it is seen because of what the seer wants, has lost, or circles. Perception of the beautiful is thus never the innocent registration of a beauty simply there, and any theory that treats it so has failed to ask who is doing the seeing and around what lack their seeing is organised. The paper will draw on this again when it comes to the reproduction of experience in silence (\xref{sec:typology}), where the psychoanalytic account of how experience is reworked after the fact becomes the theoretical body of one of its practical models.

\subsection{The objection that there is no aesthetic perception}
\label{sub:gen-dickie}

An objection now stands directly across the path, and it must be met before the argument proceeds, because if it holds there is no distinctively aesthetic perception for any education to cultivate. In his attack on the aesthetic attitude, Dickie argued that the supposed special mode of aesthetic attention is a myth: every alleged instance of a distanced or disinterested aesthetic attending turns out, on inspection, to be merely a case of attending or of failing to attend, with no special aesthetic manner of attention doing any work \citep{dickie1964}. There is attention, and there is inattention, and there is no third thing that is aesthetic attention. The paper accepts the premise and denies the conclusion, and the way it does so shows what its own claim is not. It grants Dickie that there is no special psychological state of aesthetic attention laid over ordinary perception, no distinct faculty that switches on before a sunset and off before a spreadsheet. Generative relational aesthetic perception is not the exercise of such a faculty, and a theory that posited one would deserve Dickie's refutation. But the conclusion Dickie drew, that aesthetic perception has therefore no distinctive character and reduces to ordinary attention, follows only on a further assumption that this paper rejects: that perception itself is mere registration, so that once the special attitude is removed nothing constructive remains to distinguish one perceiving from another. That assumption is exactly what the three preceding genealogies deny. If ordinary perception is already a construction of the world under a frame, then the aesthetic is not a special attitude added to a neutral perceiving but a modality of the constructive activity that all perceiving already is, a manner in which the generative work of perception is drawn upon and intensified. Dickie was right that there is no aesthetic attitude sitting on top of perception. He was wrong to conclude that perception is therefore aesthetically undifferentiated, because he took perception to be a receiving where it is a making. Removing the myth of the aesthetic attitude does not leave a flat field of mere attention; it leaves the generativity of perception itself, within which the aesthetic is a way that generativity is exercised.

\subsection{The generation of atmosphere}
\label{sub:gen-atmosphere}

The contemporary aesthetics of atmosphere completes this section by naming, in aesthetic terms proper, the relational locus that the paper's ontology requires. B\"{o}hme's new aesthetics displaces the beautiful from both of its traditional addresses. Beauty is not a property the subject projects onto a neutral object, and it is not a property the object possesses in itself awaiting discovery; it is the atmosphere, the tuned in-between in which the qualities radiating from things and the felt bodily state of the perceiver meet, and which is neither simply out there nor simply within \citep{bohme2017}. Atmosphere is, in B\"{o}hme's account, among the primary objects of perception rather than a vague addition to it, and it is generated in the between rather than found on either side. This is the aesthetic articulation of the ontological claim with which this paper began, and the paper takes B\"{o}hme as its ally in displacing beauty from subject and object alike into the relation between them. It must also mark where it goes further. B\"{o}hme's between is the between of a single perceiver and an environment; the relation he theorises runs from one felt body to the qualities of a place. The relation this paper cares about most is the further one, between two perceivers turned together toward a shared world, in which the atmosphere generated is not one body's attunement to a room but the shared generation of a perceived world by two who are also perceiving each other perceiving it. That extension is the work of the later parts of the paper. What this section has established is the ground it stands on: that across phenomenology, cognitive science, psychoanalysis, and the aesthetics of atmosphere, perception is found to generate its world under a prior frame, and beauty is found in the relational between rather than in subject or object alone. The frame's plasticity, which makes an education of perception possible, is the subject of the spine to which the paper now turns.
